
\chapter{Discussão}
\label{chap:discussao}

% Interpretar os resultados apresentados no Capítulo 15.
% Discutir em que medida eles respondem aos objetivos do trabalho
% e quais implicações têm para o desenvolvimento do DoseCerta.

\section{Análise do Atendimento aos Objetivos}
\label{sec:analise-objetivos}

% Retomar o objetivo geral e os objetivos específicos.
% Discutir quais foram alcançados, parcialmente alcançados
% ou não alcançados, com base nos resultados apresentados.

\section{Análise do Atendimento aos Requisitos}
\label{sec:analise-requisitos}

% Interpretar o grau de atendimento aos requisitos funcionais
% e não funcionais.
% Explicar as causas de requisitos não atendidos ou parcialmente
% atendidos, quando identificadas.

\section{Integração entre os Subsistemas}
\label{sec:discussao-integracao}

% Discutir como os componentes mecânico, eletrônico,
% embarcado, backend e aplicativo funcionaram em conjunto.
% Analisar dificuldades de integração e seus efeitos no sistema.

\section{Usabilidade e Adequação ao Público-Alvo}
\label{sec:discussao-usabilidade}

% Discutir a adequação das decisões de projeto às necessidades
% do público-alvo.
% Basear afirmações sobre experiência dos usuários em avaliações
% realmente realizadas. Caso não tenham ocorrido, explicitar
% que a usabilidade ainda precisa ser avaliada.

\section{Confiabilidade e Segurança}
\label{sec:discussao-confiabilidade}

% Analisar evidências relacionadas à confiabilidade e segurança
% do dispensador, considerando os testes executados.
% Explicitar riscos que não foram avaliados ou mitigados.
% Não afirmar segurança clínica com base apenas em testes técnicos.

\section{Comparação com Soluções Existentes}
\label{sec:comparacao-solucoes}

% Comparar o DoseCerta com soluções discutidas no Capítulo 3.
% Utilizar critérios consistentes, como funcionalidades,
% forma de dispensação, interação e limitações.
% Evitar alegações de superioridade sem evidências.

\section{Limitações do Projeto}
\label{sec:limitacoes-projeto}

% Descrever limitações técnicas, metodológicas e de avaliação.
% Distinguir limitações do protótipo, do processo de desenvolvimento
% e dos testes realizados.

\section{Implicações e Melhorias Futuras}
\label{sec:melhorias-futuras}

% Apontar melhorias decorrentes das limitações e resultados.
% Priorizar propostas justificadas por problemas identificados,
% requisitos pendentes ou necessidades ainda não avaliadas.

\section{Síntese da Discussão}
\label{sec:sintese-discussao}

% Consolidar a interpretação dos resultados e preparar
% a transição para as conclusões do trabalho.
